\section{An Outward Algebraic Error Account for the Scalar Comparator}
\label{app:r15scalar}

Fix the declared spatial nodes, common spacing $\Delta y$, time step $h$,
action endpoints, and terminal and boundary data.  Each stored binary64 input
denotes an exact real in this finite-grid problem.  Let
$\eta=(\sigma_I^2+\sigma_C^2)/2$,
$d_0(y)=a-\eta+\kappa\tanh y$, and define the forward and backward slopes
$p_f=(v_{i+1}-v_i)/\Delta y$ and $p_b=(v_i-v_{i-1})/\Delta y$.
The action-dependent upwind Hamiltonian is
\begin{equation}
 \mathcal H_i(m;v)=\log m-\frac\zeta2m^2
        +\max\{d_0(y_i)-m,0\}p_f+\min\{d_0(y_i)-m,0\}p_b.
 \label{eq:r15-scalar-upwind-hamiltonian}
\end{equation}
The positive-drift and negative-drift branches are each concave in $m$.
Their feasible intervals are, respectively,
$[\ell,\min\{u,d_0(y_i)\}]$ and
$[\max\{\ell,d_0(y_i)\},u]$, when nonempty.  On either branch, the
unconstrained positive critical point solves $1/m-\zeta m=p_j$.
The Howard update clips this point to the branch interval and compares the
two resulting objective values.  Thus the sign of the drift selects the
appropriate difference operator throughout maximization.

The floating-point Howard solve proposes values and actions.  A separate
outward calculation encloses the Bellman residual, including the error in
the proposed maximizer.  For a concave branch $H$ on $[a,b]$ and any
$m_\circ\in[a,b]$,
\begin{equation}
 \max_{m\in[a,b]}H(m)
 \le H(m_\circ)+
 \max\{0,H'(m_\circ)(a-m_\circ),H'(m_\circ)(b-m_\circ)\}.
 \label{eq:r15-scalar-tangent-upper}
\end{equation}
Interval evaluations of $\log$, $\tanh$, the slopes, and the arithmetic
operations enclose the right-hand side.  Interval uncertainty in $d_0(y_i)$
is handled by enlarging each branch interval before constructing its upper
bound.  An arbitrary globally feasible action supplies a lower bound for
the maximized Hamiltonian.  Combining the two bounds encloses
\begin{equation}
 r_{k,i}=
 \frac{v_{k,i}-v_{k+1,i}}h+\rho v_{k,i}
 -\eta\frac{v_{k,i+1}-2v_{k,i}+v_{k,i-1}}{(\Delta y)^2}
 -y_i-\max_m\mathcal H_i(m;v_k).
 \label{eq:r15-scalar-discrete-residual}
\end{equation}
The artifact stores every proposed value and action, the boundary arrays,
the action endpoints, and the outward residual bound at every time step.

\begin{proposition}[Algebraic error in the declared scalar boundary problem]
\label{prop:r15scalarresidual}
Fix finitely many ordered spatial coordinate labels and a common positive
difference spacing $\Delta y$,
$h>0$, $\rho\ge0$, $\eta\ge0$, compact action intervals
$[\ell_k,u_k]$ with $0<\ell_k\le u_k<\infty$, and finite terminal and
boundary data.  Let $\widetilde v$ denote the exact solution of the backward
Bellman equations~\eqref{eq:r15-scalar-discrete-residual} with zero residual.
Let $v$ have identical terminal and boundary data and residuals $r_{k,i}$.
For any $R_k\ge\max_i|r_{k,i}|$, define $E_N=0$ and
\begin{equation}
 E_k=\frac{E_{k+1}+hR_k}{1+\rho h}.
 \label{eq:r15-scalar-algebraic-bound}
\end{equation}
Then the exact finite-grid solution exists uniquely, and
$\|v_k-\widetilde v_k\|_\infty\le E_k$ at every time step.
\end{proposition}

\begin{proof}
For a fixed action $m$, write the discrete generator at an interior node as
\[
 \begin{aligned}
 (L_mw)_i&=\lambda_i^-(m)(w_{i-1}-w_i)
            +\lambda_i^+(m)(w_{i+1}-w_i),\\
 \lambda_i^-(m)&=\frac{\eta}{(\Delta y)^2}
                 +\frac{\max\{m-d_0(y_i),0\}}{\Delta y},\\
 \lambda_i^+(m)&=\frac{\eta}{(\Delta y)^2}
                 +\frac{\max\{d_0(y_i)-m,0\}}{\Delta y}.
 \end{aligned}
\]
Both rates are nonnegative.  Each frozen-policy matrix therefore has
nonpositive off-diagonal entries and diagonal dominance at least
$h^{-1}+\rho$ after the boundary values are moved to the right-hand side.
This is the $M$-matrix comparison property used below.

For completeness, existence can be established at each backward step by
uniformizing these finitely many bounded rates.  Choose
$\Lambda\ge\max_{i,m}\{\lambda_i^-(m)+\lambda_i^+(m)\}$.
When $\Lambda>0$, $L_m=\Lambda(P_m-I)$, where $P_m$ has nonnegative
entries and row sums one before the fixed boundary values are substituted.
The Bellman equation is a fixed-point equation whose dependence on the
unknown interior values has supremum-norm Lipschitz constant at most
$\Lambda/(h^{-1}+\rho+\Lambda)<1$.  The fixed-point theorem gives a
unique solution.  If $\Lambda=0$, the equation is pointwise and its
coefficient $h^{-1}+\rho$ is strictly positive.  Backward recursion starts
from the specified terminal values.

To establish the error bound, suppose
$\|v_{k+1}-\widetilde v_{k+1}\|_\infty\le E_{k+1}$ and set
$e_i=v_{k,i}-\widetilde v_{k,i}$.  If $M=\max_i e_i>0$, the maximum
occurs at an interior node $i$, because the boundary error is zero.
Nonnegative rates give $(L_me)_i\le0$ for every feasible $m$.
Writing $f_i(m)=y_i+\log m-\zeta m^2/2$, it follows that
\[
 \sup_m\{(L_mv_k)_i+f_i(m)\}
 -\sup_m\{(L_m\widetilde v_k)_i+f_i(m)\}\le0.
\]
Subtract the exact and approximate Bellman equations at this node.
The preceding inequality implies
\[
 (h^{-1}+\rho)M
 \le h^{-1}(v_{k+1,i}-\widetilde v_{k+1,i})+r_{k,i}
 \le h^{-1}E_{k+1}+R_k.
\]
Interchanging $v_k$ and $\widetilde v_k$ gives the same bound for the
largest negative error.  Equation~\eqref{eq:r15-scalar-algebraic-bound}
follows, and induction from $E_N=0$ completes the proof.
\end{proof}

The implementation evaluates this recurrence outward under the inherited
arithmetic contract.  It controls algebraic solution error for a specified
finite-grid boundary problem.  The separate grid, domain, and deployed-policy
mesh comparisons measure the remaining numerical sensitivities.

\begin{corollary}[Loss from a nested scalar action restriction]
\label{cor:r15scalartube}
Consider two problems in Proposition~\ref{prop:r15scalarresidual} with
identical spatial labels, difference and time steps, primitives, and terminal
and boundary data.  Suppose the action intervals of the second problem are
contained in those of the first at every date.  Denote their exact values by
$\widetilde v^F$ and $\widetilde v^T$, their computed values by $v^F$ and
$v^T$, and their algebraic error bounds by $E^F$ and $E^T$.  Then, at every
initial-date node,
\[
 0\le\widetilde v^F_{0,i}-\widetilde v^T_{0,i}
 \le \|v^F_0-v^T_0\|_\infty+E^F_0+E^T_0.
\]
In particular, identical stored initial-date values reduce the upper bound
to $E^F_0+E^T_0$.
\end{corollary}

\begin{proof}
At each backward step, use the common uniformization constant from the
preceding proof for both action sets.  The resulting Bellman maps are
monotone in continuation and current interior values.  Enlarging the
maximization set increases the map pointwise.  The contraction fixed points
therefore inherit this order; backward induction from the common terminal
data gives $\widetilde v^F\ge\widetilde v^T$.  The upper inequality follows
by adding and subtracting $v^F_{0,i}$ and $v^T_{0,i}$ and applying the two
algebraic error bounds.
\end{proof}
